% Generated draft: requires mathematical and source review.
\documentclass[11pt]{article}
\usepackage[margin=25mm]{geometry}
\usepackage{fontspec}
\setmainfont{Latin Modern Roman}
\usepackage{amsmath,amssymb,mathtools,unicode-math}
\setmathfont{Latin Modern Math}
\usepackage{enumitem,needspace,xurl,hyperref}
\setlength{\parindent}{0pt}
\setlength{\parskip}{5pt}
\setlength{\emergencystretch}{3em}
\newcommand{\N}{\mathbb{N}}
\newcommand{\Z}{\mathbb{Z}}
\newcommand{\Q}{\mathbb{Q}}
\newcommand{\R}{\mathbb{R}}
\newcommand{\C}{\mathbb{C}}
\newcommand{\F}{\mathbb{F}}
\newcommand{\A}{\mathbb{A}}
\newcommand{\PP}{\mathbb{P}}
\newcommand{\E}{\mathbb{E}}
\newcommand{\eps}{\varepsilon}
\newcommand{\dd}{\,\mathrm{d}}
\begin{document}
\section*{UnsolvedMath ID: 800671}
\textbf{Source problem code:} catalogue entry
\subsection{Definitions and mathematical statement}
For a topological space \(X\), use the conventional real-valued notation
\[
C_p(X)=\{f:X\to\R:f\text{ is continuous}\},
\]
with the topology of pointwise convergence. A basic neighborhood of \(f\in C_p(X)\) is
\[
[f;K,\varepsilon]=\{g\in C_p(X):|g(k)-f(k)|<\varepsilon\text{ for every }k\in K\},
\]
where \(K\subseteq X\) is finite and \(\varepsilon>0\).

A topological space \(Y\) has the \emph{Reznicenko property} if, for every \(A\subseteq Y\) and every \(y\in\overline A\setminus A\), there exist pairwise disjoint nonempty finite sets \(F_n\subseteq A\), indexed by \(n=1,2,\ldots\), such that every neighborhood of \(y\) meets all but finitely many of the sets \(F_n\).

The supplied Problem 44 asks:
\[
\text{Does }C_p(\mathbb{N})\text{ have the Reznicenko property?}
\]
The supplied extract does not define the source notation \(\mathbb{N}\) or its topology. If it means the natural numbers with the discrete topology, then \(C_p(\mathbb{N})=\R^{\mathbb{N}}\) with the product topology, and the question becomes whether the following assertion holds: for every \(A\subseteq\R^{\mathbb{N}}\) and every \(f\in\overline A\setminus A\), there are pairwise disjoint nonempty finite sets \(F_n\subseteq A\) such that
\[
\begin{split}
&\text{for every finite }K\subseteq\mathbb{N}\text{ and every }\varepsilon>0,
\text{ there exists }m\geq1\text{ such that}\\
&\qquad\text{for every }n\geq m\text{ there exists }g\in F_n
\text{ with }|g(k)-f(k)|<\varepsilon\text{ for all }k\in K.
\end{split}
\]
This specialization is conditional on the discrete-natural-number interpretation. The intended domain must be checked against the source's notation before this draft can be treated as an unambiguous transcription. No answer to the question is asserted.
\subsection{Short English statement}
The source asks whether \(C_p(\mathbb{N})\) has the Reznicenko property: at each point in the closure of a set that does not contain the point, one can find disjoint nonempty finite subsets that eventually meet every neighborhood. The meaning and topology of the source's \(\mathbb{N}\) require verification; the discrete-natural-number formulation above is conditional.
\subsection{Sources}
\begin{itemize}
\item Local UnsolvedMath registry, supplied complete record: permanent ID \(800671\), source problem code \texttt{catalogue entry}, upstream revision \texttt{d371ee917ac9647fbb06cd6fc529deec6f3dda82}.
\item Tsaban, \emph{Selection Principles in Mathematics: A Milestone of Open Problems} (2003), Problem 44. Supplied source URL: \url{https://arxiv.org/abs/math/0312182}. The record reports access on 2026-10-16 and extraction by \texttt{arxiv\_latex}; these are supplied provenance statements, not an independent source visit.
\item The supplied attribution/context is explicitly \texttt{[SakaiRez, Question 3.6]}. No full bibliographic identification of \texttt{SakaiRez} is supplied. The registry lists Tsaban as proposer and 2003 as the proposed year.
\item The registry labels this record \texttt{open}, with update date 2026-10-16. This is supplied status metadata, not an independently verified literature assessment.
\end{itemize}
\end{document}
